\chapter{Procedimento para verificação ou demonstração prática}
\label{cap:demonstracao}

A demonstração deve tornar visível a ligação entre perceção, inferência e ação.
O roteiro seguinte cobre os critérios de base de conhecimento, motor,
explicações, metaconhecimento e produto final.

\begin{enumerate}
  \item Iniciar o \emph{backend} e o \emph{frontend} usando o \emph{script}
    adequado ao sistema operativo do utilizador na pasta \cod{scripts}, e
    selecionar o motor \enquote{Prolog} na interface Web.
  \item Premir \enquote{Start}.
  \item Carregar um cenário sem perigos e mostrar a decisão base de manutenção da
    velocidade.
  \item Introduzir um peão na trajetória e observar o evento, a regra disparada,
    a proposta e a travagem.
  \item Abrir a memória de trabalho e localizar \cod{pedestrian},
    \cod{pedestrian_conflict} e \cod{proposal}.
  \item Mostrar a evidência da proposta para responder como e por que a decisão
    foi obtida.
  \item Criar um obstáculo com uma faixa adjacente ocupada e mostrar os
    subobjetivos que explicam por que não ocorreu a mudança.
  \item Executar três amostras TTC críticas dentro de dois segundos e confirmar a
    regra temporal de emergência.
  \item Alterar visibilidade e aderência para demonstrar fuzificação, conclusão
    difusa e teto de velocidade.
  \item Trocar para Drools no mesmo cenário e comparar ação, velocidade e regra.
  \item Apresentar as métricas \cod{rematches}, \cod{firings} e \cod{ruleMillis}
    do último ciclo.
\end{enumerate}
